\section{结论与展望}

\subsection{主要贡献}

本文面向F1赛车气动稳定性优化问题，
构建了一套完整的数据驱动代理优化框架。
研究的过程本身揭示了一个比PSO寻优性能更重要的发现：
代理模型驱动优化的效果上限并非由优化器决定，
而是由数据结构和目标函数设计共同约束。
以下从五个技术贡献和一个核心科学发现两个层面总结本研究的成果。

\vspace{0.2cm}

\begin{center}
\fbox{\begin{minipage}{0.92\textwidth}
\vspace{0.2cm}
\textbf{本研究最重要的科学发现}

景观诊断揭示了代理模型驱动优化中的一个方法论核心结论：
当训练数据中86\%的样本stability$\geq 99$时，
单目标优化景观趋于平坦——在此条件下，
PSO、GA、DE的寻优结果几乎无差异（极差仅0.15）。
这一发现表明，在代理模型驱动的优化中，
\textbf{优化算法并非瓶颈——数据分布结构才是优化效果的根本约束}。
该结论将本研究的焦点从"哪种算法更优"提升到
"数据驱动的优化方法何时有意义"这一更深层次的问题。
\vspace{0.1cm}
\end{minipage}}
\end{center}

\vspace{0.3cm}

\noindent
\textbf{（1）代理模型谱系对比与优选}

在6个模型（Ridge$\to$DeepEnsemble）的完整谱系中进行了系统的精度-效率评估，
选定XGBoost（R$^2$=0.9999, MSE=0.0312, 延迟0.242ms）
作为可解释性分析模型，
DeepEnsemble作为PSO优化模型（R$^2$=0.9995, 含不确定性估计）。
全部模型的R$^2$超过0.92理想目标。

\vspace{0.3cm}

\noindent
\textbf{（2）不确定性感知的风险敏感优化（重要方法设计）}

提出$\mu(x) - \lambda\sigma(x)$风险敏感适应度函数，
利用DeepEnsemble的不确定性估计同时解决
OOD幻觉（A1消融证实：移除$\sigma$后适应度飙升至101.30）、
数据不平衡偏置和代理模型可靠性三个问题。
消融实验系统验证了$\sigma$惩罚作为OOD防护屏障的不可替代性。

\vspace{0.3cm}

\noindent
\textbf{（3）不确定性引导的模型精炼}

实现了PSO探索反馈$\to$模型精炼的闭环机制。
一轮精炼即可将MSE从0.3272降至0.1444（降低55.9\%），
PSO收敛速度提升2.1倍。

\vspace{0.3cm}

\noindent
\textbf{（4）多场景加权多目标优化+可解释性体系}

在4组F1赛道工况下完成了多目标PSO寻优，
引入气动效率和功率消耗作为竞争性目标，
使场景间区分度从0.1\%提升至23\%以上。
建立了SHAP+排列重要性+约束灵敏度+Landscape Diagnosis
的可解释性验证体系，
通过多种独立方法交叉验证了分析结论的可信度。

\vspace{0.3cm}

\noindent
\textbf{（5）Porpoising风险空间刻画}

基于DeepEnsemble梯度构建了一阶风险热力图
（高速+小翼角区域风险为低速+大翼角的6.5倍），
并通过Hessian曲率分析揭示了MLP预测曲面
在高稳定性区的极度平坦特性。

\subsection{量化指标达成情况}

\begin{table}[H]
\centering
\caption{量化指标汇总}
\begin{tabular}{lccc}
\toprule
指标 & 实际值 & 最低达标 & 理想目标 \\
\midrule
代理模型 R$^2$ (XGBoost) & \textbf{0.9999} & 0.85 & 0.92 \\
代理模型 MSE (XGBoost) & \textbf{0.0312} & 0.05 & 0.02 \\
PSO 收敛代数 (标准) & \textbf{27.9} & 80 & 50 \\
PSO 收敛代数 (自适应) & \textbf{22.1} & 80 & 50 \\
最优解标准差 (20次) & 0.0836 & 0.05 & 0.02 \\
单次寻优耗时 & \textbf{83ms} & 20s & 5s \\
\bottomrule
\end{tabular}
\label{tab:quantitative}
\end{table}

表\ref{tab:quantitative}显示：
R$^2$、PSO收敛代数、单次寻优耗时三项指标
同时达到了最低达标和理想目标；
最优解标准差接近达标（0.0836 vs 0.05），
主要受DeepEnsemble在饱和区域的输出一致性和$\sigma$惩罚项影响，
在分析中已作为研究发现的深度加以利用。

\subsection{不足与讨论}

本研究存在以下局限性：

\begin{enumerate}
    \item \textbf{数据驱动的物理极限}：
        数据集86\%的样本stability$\geq 99$，
        导致代理模型在高稳定性区输出饱和。
        所有优化器（无论PSO、梯度下降或替代方案）
        都将趋于相同的边界解。
        高比例稳定样本（86\%）导致目标函数退化，
使单目标优化景观趋于平坦——在此条件下，
PSO、GA、DE的寻优结果几乎无差异（极差仅0.15）。
该发现表明，在代理模型驱动优化中，
数据质量和目标函数设计往往比优化器选择更为关键。
这一结论不应被视为实验失败，
而是对数据驱动优化方法论边界条件的实证刻画。
    \item \textbf{代理模型的单向性}：
        真实物理中下压力对稳定性有正向贡献
        （更大的下压力提升过弯抓地力），
        但当前数据集将高下压力与低稳定性关联
        （r = $-$0.773），
        可能是赛道条件或数据采集过程的混杂因素；
    \item \textbf{模型精炼的深度}：
        仅完成了一轮浅层精炼，
        未探索多轮迭代的收敛行为和边际收益递减特性；
    \item \textbf{海豚跳的间接性}：
        Porpoising风险分析基于代理模型的梯度，
        属于间接估计。
        更直接的方法需要CFD或实验数据的支撑。
\end{enumerate}

\subsection{未来工作}

基于本研究的发现和不足，
建议从以下方向继续推进：

\begin{enumerate}
    \item \textbf{物理信息代理模型}：
        将气动平衡方程作为物理约束嵌入神经网络训练过程，
        有望打破当前数据不平衡导致的模型饱和瓶颈；
    \item \textbf{多轮主动学习}：
        探索迭代精炼的边际收益递减特性，
        确定最优的主动学习轮次；
    \item \textbf{CFD-代理混合策略}：
        在主动学习框架中，
        以限量的CFD真实仿真替代简单的训练集样本检索，
        实现真正的"仿真-代理"协同优化；
    \item \textbf{更丰富的多目标维度}：
        引入俯仰灵敏度、底盘振动频率等
        与海豚跳直接相关的物理维度，
        构建面向真实工程需求的完整多目标优化框架；
    \item \textbf{贝叶斯优化对比}：
        将风险敏感PSO与贝叶斯优化
        （利用DeepEnsemble的$\sigma(x)$做采集函数）
        进行系统对比，探索不确定性量化在两种优化范式中的差异性表现。
\end{enumerate}
